%% table 객체: tab:terms — 게임·자료 용어 (docs/tog_manuscript/sec_background.tex Table tab:terms의 국문 요약; 출처는 각 행의 bib 키)
\begin{table}[!t]
  \centering
  \caption{게임과 자료 용어}
  \label{tab:terms}
  \footnotesize
  \begin{tabularx}{\linewidth}{@{}>{\raggedright\arraybackslash}p{3.0cm}X>{\raggedright\arraybackslash}p{3.4cm}@{}}
    \toprule
    용어 & 뜻 & 출처 \\
    \midrule
    소환사의 협곡 & 다섯 명씩 두 팀이 겨루는 대칭 지도; 세 라인과 정글, 강 & \cite{lolwiki2025summonersrift,riotapi2024matchv5} \\
    챔피언 & 플레이어가 조종하는 캐릭터; 팀당 다섯 & \cite{lolwiki2026champion} \\
    골드·경험치 & 미니언·몬스터·킬에서 얻는 자원; 아이템과 레벨로 전환되어 챔피언을 강하게 함 & \cite{lolwiki2026gold,lolwiki2026experience} \\
    킬·어시스트·사망 & 킬은 킬러에게 골드를 줌; 어시스트는 직전 기여자에게 인정; 사망한 챔피언은 일정 시간 경기에서 제외 & \cite{lolwiki2026kill,lolwiki2026assist,lolwiki2026death} \\
    포탑·억제기·넥서스 & 라인을 지키는 구조물; 넥서스 파괴가 승리 조건 & \cite{lolwiki2026turret,lolwiki2026inhibitor,lolwiki2026nexus} \\
    엘리트 몬스터 & 드래곤, 협곡의 전령, 공허 유충, 아타칸, 바론; 처치한 팀에 보상 & \cite{lolwiki2026dragonpit,lolwiki2026riftherald,lolwiki2026voidgrub,lolwiki2026atakhan,lolwiki2026baron} \\
    시야·와드 & 시야 밖의 적 챔피언은 보이지 않음; 와드는 시야를 제공 & \cite{lolwiki2026sight,lolwiki2025ward} \\
    패치 & Riot이 규칙을 바꾸는 단위; Match-V5는 15.N, 공개 번호는 25.N & \cite{riot2025patch25s11,lolwiki2026v2514} \\
    참가자 프레임 & Match-V5 timeline의 약 60초 간격 상태 기록(열 명의 위치·레벨·골드·체력) & \cite{riotapi2024timeline} \\
    사건 & Match-V5 timeline의 밀리초 시각 기록(킬·구조물·몬스터 등); 챔피언 킬만 시각과 위치를 모두 가짐 & \cite{riotapi2024timeline} \\
    교전 (본 논문) & 한타와 소규모 교전을 아우르는 분석 단위. 킬을 기준으로 시간·공간 범위와 참여 조건을 정해 구성한다 & 제 \ref{ch:engagement} 장 \\
    한타 & 교전 가운데 적게 모인 쪽이 네 명 이상인 것. 플레이어는 인원이 많이 모인 맞붙음을 이렇게 부른다 & \cite{lolwiki2026terminology}, 제 \ref{ch:engagement} 장 \\
    소규모 교전 & 교전 가운데 적게 모인 쪽이 두 명 또는 세 명인 것 & 제 \ref{ch:engagement} 장 \\
    \bottomrule
  \end{tabularx}
\end{table}
